\chapter{优化模型、数值实验与案例分析}

当前项目中最有研究价值的部分，并不是某一条公式本身，而是这些公式共同构成了一个高度离散、带多重上限、且充满取整跳点的优化问题。本章在前文机制整理的基础上，把宠物培养与战斗目标统一写成可计算问题，并讨论当前项目为什么选择“先代理优化，再回到完整战斗模型验证”的技术路线。

\section{优化问题的形式化}

把整个宠物培养过程写成统一映射后，可以把最终优化问题抽象为
\begin{equation}
    \max_{x \in \mathcal{X}} U\bigl(\Psi(x)\bigr),
    \label{eq:global-opt}
\end{equation}
其中 $x$ 表示决策变量，$\mathcal{X}$ 表示所有可行培养方案，$\Psi$ 表示“培养方案 $\rightarrow$ 成长档 $\rightarrow$ 随机成长 $\rightarrow$ 战斗属性 $\rightarrow$ 战斗表现”的复合映射，$U$ 则是研究者真正关心的目标函数。

在当前工具中，$x$ 至少可以包含三类变量：转生输入、融合输入和喂药方案。它们之中，喂药方案最适合先被形式化为整数规划。若记第 $i$ 维第 $\ell$ 级药的颗数为 $c_{i\ell}$，则其约束满足
\begin{equation}
    c_{i\ell} \in \mathbb{Z}_{\ge 0},
    \qquad
    \sum_{i=1}^{4}\sum_{\ell=1}^{5} c_{i\ell} = 40.
    \label{eq:feed-constraint}
\end{equation}
但由于后续还存在折算、单项上限、总和压缩和比例回分配，式~\eqref{eq:feed-constraint} 只是最外层约束，而不是完整可行域描述。

\section{蛋喂养最优化}

在喂药问题中，最自然的目标是直接最大化某个最终战斗指标，但这会把离散喂药、孵化、出生随机、升级随机和战斗公式全部卷到同一个搜索里，复杂度极高。因此，当前项目采用了更务实的两层方法：先在孵化后成长 $\bm{h}$ 上做离散优化，再把候选解带回更完整的战斗模型验证。设蛋基础成长为 $\bm{b}$，喂药方案为 $c$，则喂药层的中间问题可写为
\begin{equation}
    \max_{c} U_0\bigl(\bm{h}(c;\bm{b})\bigr),
    \label{eq:hatch-opt}
\end{equation}
其中 $U_0$ 可以是线性目标，也可以是词典序目标。

这一层之所以必须显式建模，是因为“全部喂 5 级药”并不总是最优。若药效在折算后已经触发 50 压缩或最终 150 压缩，那么继续把药往同一维堆高，可能只会把别的维度一起拉进压缩区，而不再提高目标维度的最终收益。也正因为此，在某些高基础蛋上，混入较低等级药反而可能比全 5 级药更优。这个现象不是经验技巧，而是由多次向下取整和总和上限共同造成的离散最优化结果。

当前项目还引入了词典序目标，即先最大化主目标，再在主目标并列解中最大化次目标。这一步的必要性来自大量退化解。例如，在药效原始值满足线性关系的情况下，不同药级组合可能给出完全相同的折算结果。若没有次级目标或其它偏好规则，输出候选解就会极度冗余。当前实现里偏向更高等级素材，本质上就是在这些退化解之间加入了一个稳定的优先级。

\section{战斗目标反推}

真正的战斗目标往往不是“让某个孵化后成长最大”，而是“让先手概率、命中率、会心率或骑宠伤害最大”。因此，需要把战斗目标反投影回成长空间。若记战斗目标为
\begin{equation}
    U_{\mathrm{battle}} = U\bigl(HP,ATK,DEF,SPD,p_{\mathrm{hit}},p_{\mathrm{crit}},damage,\dots\bigr),
\end{equation}
则一个常见策略是先在当前点做局部线性化，把它近似为
\begin{equation}
    U_{\mathrm{battle}} \approx w_V h_V + w_S h_S + w_T h_T + w_D h_D + const,
\end{equation}
再回到式~\eqref{eq:hatch-opt} 求解。

例如，若目标是提升玩家对玩家的先手与会心，则由于这两项都对敏捷高度敏感，线性化后通常会给 $w_D$ 更高权重；若目标是骑宠近战爆发，则 $w_S$ 与 $w_D$ 往往都会较高，因为 $S$ 主要影响攻击而 $D$ 同时影响先手和会心；若目标是承伤，则应提高 $w_V$ 与 $w_T$。这一“战斗目标 $\rightarrow$ 成长代理目标”的反推步骤，正是工具中各种喂药搜索模式背后的统一数学结构。若后续把装备、阵型、AI 行为与地图场地一并纳入模型，那么这里的局部线性化对象也不必只停留在 $HP,ATK,DEF,SPD$ 与少数概率量上，而可以扩展成“培养方案 $\rightarrow$ 完整战斗决策收益”的多层代理目标。换言之，当前的成长代理优化并不是一个封闭终点，而是后续做交互式战斗反推、自动配装和全链路求解的中间层。

\section{转生与融合策略比较}

从优化视角看，转生与融合并不是同层决策。转生更像“先重塑成长档，再把它交给完整成长随机”；融合更像“把多只宠物的成长信息压缩进一个蛋面中间态”。因此，两者各自的瓶颈也不同。

转生的主要瓶颈来自转生能力上限。无论 MM、等级和原宠总和如何继续提高，转生能力值最终都会被 $A_{\max}$ 截断，因此高端区间往往更像“方向重分配问题”，而不是“总和继续增长问题”。融合的主要瓶颈则来自逐层取整：低等级宠先做 $0.8$ 折扣，再参与主副权重分配，最后蛋又要经过孵化压缩。这意味着融合的损失是沿途积累的，而不是最后一步才一次性发生。

由此可以得到一个很实用的比较框架：若当前手里已有高总和、高方向一致性的成熟宠，转生更适合作为“保留方向并整体抬高”的策略；若目标是把多只宠物的局部优势重新拼装到一只新宠上，融合与喂药则更适合作为“重构方向”的策略。但只要问题进入压缩区，无论是转生还是融合，都不再适合用“多一点总和一定更好”的线性直觉来判断。

\section{案例研究}

下面给出三个不依赖具体服务器数值、但足以说明机制的案例。

\subsection{案例一：先手优先的骑宠近战}

设目标是让一只骑宠近战配置在 PVP 中尽量抢到先手，同时保留一定爆发。由第 6 章和第 10 章可知，骑宠近战时速度参照与攻击都更偏宠物，因此代理目标应更偏向 $D$ 与 $S$。一种自然的近似目标是
\begin{equation}
    U_1 = \alpha h_D + \beta h_S,
    \qquad \alpha > \beta > 0.
\end{equation}
若基础蛋本来就有较高总和，则继续把所有高等级药堆到 $D$ 上，可能因为触发总和压缩而连带损失 $S$，从而在最终骑宠伤害上得不偿失。这正是词典序目标或多候选解比较有价值的场景。

\subsection{案例二：高基础蛋的压缩陷阱}

设一枚蛋的基础成长总和已经接近最终上限。若继续用“每一颗药都尽量给最高级”的策略，则折算后的工作值很容易触发 50 压缩，而孵化后成长又可能触发 150 压缩。此时，药效的边际收益不再等于原始药值。对这类问题，最优方案常常不是单纯堆单项，而是控制进入压缩前的分布，使目标维度在压缩后仍保留最大整数份额。也就是说，离散最优解更像“卡阈值”而不是“无限堆高”。

\subsection{案例三：主目标并列下的次目标选择}

设若干喂药方案都能把主目标维度推到同一上限，此时单看主目标已经无法区分优劣。若没有次目标或偏好规则，输出会出现大量本质等价的退化解。当前项目采用“主目标 + 次目标 + 高等级素材优先”的方案，本质上是在一个大并列集合里引入稳定排序。这样做的意义不是改变底层规则，而是让结果更适合人类使用和二次微调。

综合这些案例可以看到，当前项目中的优化问题并不是“算力不够才退而求其次”的粗糙近似，而是因为底层机制本身就是一个离散、分段、带上限的系统。与其追求一个表面上更统一但实际上更难验证的全局闭式解，不如先沿着源码结构把问题拆成可验证的中间层，再逐层向战斗目标推进。这也是整篇论文所采用的方法论。对工具实现而言，这意味着前端并不需要一次性吃下“完整战斗模拟”这一最重问题，而可以先以孵化后成长、再以高等级面板、最后以战斗代理收益三层逐步接入；对后续研究而言，这也为更系统的蒙特卡洛实验、PVP/PVE 参数扫描以及多目标求解保留了自然的扩展接口。

